\section{Related Work}
\label{sec:related}

\textbf{Parameter-Space Model Merging.} The field of model merging seeks to
combine multiple neural networks with identical architectures into a single
model. Early works focused on uniform weight averaging \cite{fedavg,
wortsman2022model}. Task Arithmetic \cite{ilharco2022editing} showed that
editing pre-trained models via linear combinations of task vectors can
successfully add or remove capabilities. TIES-Merging \cite{yadav2023ties}
resolved parameter interference by pruning redundant signs and resolving sign
conflicts before averaging. RegCalMerge \cite{trial2_submission1} addressed
"sacrificial task bias" via class-capacity normalization, while PolyMerge
\cite{trial2_submission3} restricted coefficients to low-degree polynomials. Our
work differs by deriving the spatial relationship of layer coefficients from a
first-principles PAC-Bayes framework.

\textbf{Test-Time Adaptation and Calibration.} Test-Time Adaptation (TTA) adapts
a model to a target distribution during inference using only unlabeled data.
TENT \cite{tent} minimizes prediction entropy on incoming test streams. In the
context of model merging, AdaMerging \cite{yang2024adamerging} optimizes
layer-wise coefficients at test time using entropy minimization. However,
unconstrained TTA is prone to representation collapse under small batch sizes or
covariate shifts. In this work, we expose the \emph{Overfitting-Optimizer
Paradox} in test-time merging and show how unconstrained optimization causes
catastrophic generalization collapse, which we resolve through rigorous
PAC-Bayes regularization.

\textbf{PAC-Bayes Generalization and Spatial Priors.} PAC-Bayes theory provides
mathematically rigorous generalization guarantees for randomized classifiers
\cite{mcallester1999some, catoni2007pac}. It bounds the generalization error of
a posterior distribution $Q$ using its empirical error and the KL divergence to
a prior distribution $P$. While traditionally used to derive weight-space
generalization bounds \cite{dziugaite2017computing}, we apply PAC-Bayes theory
directly to the space of test-time merging coefficients. By modeling the prior
$P$ as a Gaussian Process (GP) over normalized network depth, we capture the
physical, spatial correlation of neural representation layers. This formulation
leverages the rich literature on GP priors in deep learning
\cite{neal1996priors, rasmussen2006gaussian} to establish a mathematically
sound, non-heuristic spatial regularizer.
